\documentclass[../../main.tex]{subfiles}
\begin{document}

\begin{center}
  {\fontsize{48pt}{58pt}\selectfont \textbf{1983}\normalsize\textbf{年度}}
\end{center}

\begin{table}[htb]
  \centering
  \caption{出題テーマと難易度}
  \begin{tabular}{|c|p{9cm}|c|}
    \hline
    \textbf{問題番号} & \textbf{テーマ} & \textbf{難易度} \\
    \hline
    第1問 & 3校の連勝制野球大会での優勝確率 & \\
    \hline
    第2問 & 直角三角形に付けた台形の形 & \\
    \hline
    第3問 & 行列の回転・拡大分解と一次変換 & \\
    \hline
    第4問 & 直交する3辺に下ろした垂線と体積比 & \\
    \hline
    第5問 & 動く2点が出会うまでの時間の最小化 & \\
    \hline
    第6問 & 円錐容器からの排水と水面の高さ & \\
    \hline
  \end{tabular}
\end{table}

\begin{table}[htb]
  \centering
  \caption{解答}
  \renewcommand{\arraystretch}{1.4}
  \begin{tabular}{|c|c|p{9cm}|}
    \hline
    \textbf{問題番号} & \textbf{小問} & \textbf{答え} \\
    \hline
    第1問 & (1) & $\dfrac{2pq}{2-p(1-q)}$ \\
    \cline{2-3}
    & (2) & $\dfrac{pq(1-p)}{2-q(1-p)}$ \\
    \cline{2-3}
    & (3) & $pq\left(\dfrac{1}{2-p(1-q)}+\dfrac{1}{2-q(1-p)}\right)$ \\
    \hline
    第2問 & (1) & $Q,A,R$は一直線上にある \\
    \cline{2-3}
    & (2) & 台形$BCRQ$は長方形 \\
    \hline
    第3問 & -- &  \\
    \hline
    第4問 & (1) & $\dfrac{V_{OPQR}}{V_{OABC}}=\dfrac{2a^2b^2c^2}{(a^2+b^2)(b^2+c^2)(c^2+a^2)}$ \\
    \cline{2-3}
    & (2) & 体積比は$\dfrac{1}{4}$以下 \\
    \hline
    第5問 & -- & $t=\dfrac{2a+b}{3}$ \\
    \hline
    第6問 & -- & $x(t)=h\sqrt{1-\dfrac{2Skt}{\pi R^2}}$ \\
    \hline
  \end{tabular}
\end{table}

\paragraph{試験形態}
6問．試験時間は要確認．

\end{document}
